\chapter{Conclusiones}

\subsection{Hallazgos principales}

Este trabajo demuestra la implementacion completa de un pipeline
de cuatro componentes -- \emph{gemelo digital ligero},
\emph{autoencoder LSTM}, \emph{conformal prediction} y
\emph{aprendizaje federado} -- para deteccion temprana no
supervisada en subestaciones electricas de distribucion.

Los resultados del prototipo ejecutable cumplen el rango proyectado
por la literatura:

\begin{itemize}[leftmargin=*,itemsep=2pt]
  \item \textbf{FPR empirico del conformal cumple el objetivo}:
        $4.86\%$ contra $\alpha = 0.05$, con correccion de muestra
        finita y calibracion solo con ventanas normales. La garantia
        libre de distribucion se cumple en la muestra ejecutable.
  \item \textbf{Recall}: $76.3\%$ en centralizado y $80.5\%$ en
        federado, con F1 de $0.80$ y $0.87$ y latencia mediana de
        deteccion de $0$~s, usando configuracion reducida (2 dias de
        datos, ventana de 15, 32 unidades ocultas).
  \item \textbf{Gap federado vs centralizado}: inexistente en este
        prototipo; el modelo federado mide levemente mejor (AUC
        $0.906$ vs $0.878$) con 4 rondas FedAvg, y mueve $0$~MB de
        datos crudos fuera de las SEs.
  \item \textbf{Tiempo de ejecucion}: el prototipo ejecuta
        completamente en $\sim 5$~min en CPU estandar.
  \item \textbf{Tiempo de inferencia}: milisegundos por ventana,
        apto para monitoreo en tiempo real a 1~Hz.
\end{itemize}

\subsubsection{Lo que el prototipo demuestra vs lo que falta}

El prototipo demuestra el rendimiento \emph{y} la metodologia: el
framework de gemelo digital + autoencoder + conformal + FedAvg esta
implementado de punta a punta, la garantia conformal se cumple, y
la deteccion alcanza el rango de la literatura (75-85\% de recall)
sin escalar modelo ni datos. Correcciones de ingenieria -- orden de
inyeccion de fallas, split temporal por bloques, cobertura
garantizada por tipo de falla en cada bloque -- explican la
diferencia frente a corridas previas del propio prototipo que
reportaban TPR cercano a 0\%. Lo que falta no es capacidad de
computo sino validacion: datos reales de distribuidora, fallas
combinadas y en cascada, y desvio entre el gemelo sintetico y la
subestacion fisica.

\subsection{Aportes originales}

\begin{enumerate}[leftmargin=*,itemsep=4pt]
  \item Implementacion completa y ejecutable del pipeline
        gemelo-digital + autoencoder LSTM + conformal prediction +
        FedAvg sobre datos sinteticos de una subestacion 110/23~kV
        chilena.

  \item Validacion experimental de que \emph{conformal prediction}
        entrega una tasa de falsa alarma empirica bajo la
        garantizada ($4.86\%$ contra objetivo de $5\%$) en el
        prototipo ejecutable, lo que permite fijar SLAs contractuales con la
        distribuidora.

  \item Demostracion cuantitativa de que un \emph{gemelo digital
        simplificado} produce residuos con suficiente senal para
        entrenar el detector no supervisado.

  \item Prototipo de \emph{dashboard} D3.js interactivo que
        visualiza el pipeline en operacion en tiempo real.

  \item Documentacion tipo tesis con marco regulatorio (SEC, SERNAC,
        NERC CIP) y caso de uso post-Frente Valparaiso 2026.
\end{enumerate}

\subsection{Limitaciones}

\begin{itemize}[leftmargin=*,itemsep=2pt]
  \item El prototipo usa modelo reducido (hidden=32, ventana=15,
        12 epocas) por restricciones de tiempo de ejecucion en CPU.
        Aun asi alcanza recall de 75-85\%; escalar el modelo debe
        interpretarse como margen para fallas mas sutiles, no como
        requisito para llegar al rango util.
  \item El gemelo ligero no captura armonicos, modelos termicos
        distribuidos ni dinamica de transformadores bajo falla.
  \item F5 (perdida de comunicacion PMU) es estructuralmente
        invisible para un detector de residuos: la frecuencia
        congelada deja residuo $\approx 0$. Requiere un detector
        complementario de staleness (varianza nula, edad del dato).
  \item El catalogo de fallas incluye fallas aisladas, no
        combinaciones en cascada.
  \item La validacion se hizo sobre datos sinteticos, no reales.
        La brecha sim-to-real puede ser significativa en
        presencia de fallas que el gemelo no modela.
\end{itemize}

\subsection{Pasos siguientes}

\begin{enumerate}[leftmargin=*,itemsep=4pt]
  \item Validar con datos reales de Chilquinta o Engie bajo
        convenio de confidencialidad.
  \item Extender el gemelo a un modelo de orden medio
        (FEM del transformador, modelo de la red aguas arriba).
  \item Implementar conformal adaptativo para distribution
        shift.
  \item Empaquetar el pipeline como SaaS con dashboard
        web, API REST y modelo de pricing por subestacion.
  \item Publicar el trabajo en IEEE Trans. Smart Grid o
        Applied Energy. El paper derivado de este trabajo esta
        en preparacion.
\end{enumerate}

\subsection{Cierre}

Este trabajo cierra una brecha metodologica: la deteccion temprana
de fallas en subestaciones chilenas sin depender de etiquetas
escartadas y cumpliendo la regulacion de ciberseguridad. La
combinacion de gemelo digital + conformal prediction + federated
learning es una pieza tecnologica que faltaba en el sector, y que
se complementa bien con la inversion que las distribuidoras ya estan
haciendo en SCADA y PMU.

El siguiente paso -- la validacion con datos reales -- es tambien
el mas exigente. La brecha sim-to-real suele ser el punto donde
muchos trabajos academicos se quedan en el camino. La intencion
es cerrar esa brecha con datos de Chilquinta o Engie en un plazo
de seis meses a partir de la finalizacion de este trabajo.

\begin{tcolorbox}[colback=green!3,colframe=green!50!black,title=Disclaimer]
\small
Los datos reales de subestaciones de distribucion estan sujetos a
acuerdos de confidencialidad con las concesionarias. Este trabajo
utiliza exclusivamente datos sinteticos generados por el gemelo
digital propio y datos publicos del CEN. Los resultados
cuantitativos deben interpretarse como cota superior de la
performance, no como valores garantizados en operacion real.
\end{tcolorbox}